\section{Resultados Obtenidos}

Luego de la ejecución de la metaheurística \textit{Simulated Annealing}, se logro determinar una solución que cumpla con las demandas solicitadas por las estaciones de servicio a su vez que respete las restricciones dadas, a lo largo de la ejecución se fue transformando la solución inicial usando movimientos vecinos para lograr llegar a la solucion final.

\subsection{Configuración Experimental}

Las pruebas fueron realizadas utilizando los siguientes parámetros:

\begin{itemize}
\item Temperatura inicial: $T_0 = 1000$.
\item Temperatura final: $T_f = 0.01$.
\item Factor de enfriamiento: $\alpha = 0.95$.
\item Iteraciones por temperatura: $100$.
\item Semilla aleatoria: $64$.
\end{itemize}

La solución inicial fue generada mediante el procedimiento constructivo descrito anteriormente y posteriormente mejorada a través de los movimientos vecinos definidos para la exploración del espacio de búsqueda.

\subsection{Mejor Solución Encontrada}

La mejor solución obtenida por la metaheurística presentó un valor de función objetivo igual a:

\begin{equation}
f(S^*) = 230 + 900 + 0 + 0 = 1130
\end{equation}

donde:

\begin{itemize}
\item Costo por distancia: $C_{dist}(S^*) = 230$.
\item Costos fijos de utilización de camiones: $C_{fijo}(S^*) = 900$.
\item Penalización por demanda no satisfecha: $C_{shortage}(S^*) = 0$.
\item Penalización por restricciones: $P_{inf}(S^*) = 0$.
\end{itemize}

La distribución obtenida para los camiones fue la siguiente:

\[
R_{T1} = [3]
\qquad
R_{T2} = [4,2,1]
\]

En esta solución, el camión $T1$ abastece la estación 3, mientras que el camión $T2$ abastece las estaciones 4, 2 y 1.

La asignación de combustibles a los compartimentos fue:

\begin{table}[H]
\centering
\caption{Asignación de productos a compartimentos}
\label{tab:compartimentos}
\begin{tabular}{c|cc}
\toprule
Camión & Compartimento & Producto \\
\midrule
T1 & C0 & Regular \\
T1 & C1 & Diésel \\
T2 & C0 & DIésel \\
T2 & C1 & Regular \\
\bottomrule
\end{tabular}
\end{table}

Asimismo, las cantidades entregadas a cada estación se resumen en la Tabla \ref{tab:entregas}.

\begin{table}[H]
\centering
\caption{Cantidades entregadas por estación}
\label{tab:entregas}
\begin{tabular}{c|ccc}
\toprule
Estación & Regular (L) & Diésel (L) & Shortage (L) \\
\midrule
1 & 3000 & 2000 & 0 \\
2 & 4000 & 1500 & 0 \\
3 & 2500 & 3000 & 0 \\
4 & 1000 & 2500 & 0 \\
\bottomrule
\end{tabular}
\end{table}

\subsection{Validación de Restricciones}

La mejor solución obtenida fue sometida al proceso de validación descrito previamente. Los resultados indicaron que:

\begin{itemize}
\item Todas las rutas comienzan y terminan en el depósito.
\item Ninguna estación fue atendida por más de un camión.
\item No se excedieron las capacidades máximas de los compartimentos.
\item Cada compartimento transportó un único tipo de combustible.
\item Se respetó el límite máximo de diferencia entre \textit{fill ratios}.
\item Las ventanas de tiempo fueron satisfechas.
\item La demanda fue completamente satisfecha (o el shortage fue correctamente contabilizado).
\end{itemize}

Por lo tanto, la solución obtenida puede considerarse operacionalmente factible y representa la mejor alternativa encontrada por la metaheurística dentro de las condiciones evaluadas.

\subsection{Análisis de Resultados}

Los resultados obtenidos muestran que la metaheurística de \textit{Simulated Annealing} logra otorgar soluciones de calidad que responde a las demandas de combustible respetando restricciones, esto dado a que pudo encontrar la combinación de asignación de compartimientos, carga inicial y rutas para los camiones que tuviese el menor costo operacional posible.

La factibilidad de la solución se logro tras comprobar que esta cumpliese con las restricciones del problema de distribución de combustible, logrando un equilibrio entre la minimización de costes y cumplimiento de restricciones.

En conclusión, la metaheurística propuesta logra ser una alternativa adecuada para revolver el problema cuando los metodos exactos resultan computacionalmente complejos en diversas instancias del problema.
